% preprint/modules/background.tex
% Shared-Horizon Preprint — Narrative Modules
% Module: Background and Related Work

\section{Background and Related Work}
\label{sec:background}

Artificial intelligence research has diversified into multiple conceptual
lineages, each grounded in distinct mathematical, expressive, or structural
assumptions. This diversification has produced powerful but incompatible
frameworks for understanding system behavior. Without a neutral comparative
structure, attempts to analyze or relate these lineages often result in
conceptual collapse, where the invariants of one lineage are incorrectly mapped
onto another.

\subsection{Fragmentation of AI Lineages}

The fragmentation can be broadly categorized into five major lineages:

\begin{enumerate}
    \item \textbf{Geometric Operator Systems (NDH)}  
    Systems defined by reversible operators, commutative geometry, and
    holonomy-preserving transformations. These systems emphasize structural
    invariants and geometric stability.

    \item \textbf{Expressive Phenomenological Systems (ANIMA)}  
    Models that operate on expressive curvature, affective gradients, and
    phenomenological coherence. These systems prioritize expressive manifolds
    over operator geometry.

    \item \textbf{Pre-Classical Continuity Systems (PRECL)}  
    Frameworks that emphasize smooth continuity, pre-classical transitions, and
    low-curvature expressive surfaces.

    \item \textbf{Reversible Narrative Systems (RNE)}  
    Systems defined by reversible narrative operators, curvature-preserving
    transitions, and narrative geodesics. This lineage corrects earlier
    misclassification as symbolic ambiguity.

    \item \textbf{Structural-AI Spectral Systems}  
    Architectures defined by spectral-runtime operators, bounded spectral drift,
    and operator-driven stability envelopes.
\end{enumerate}

Each lineage has its own mathematical invariants, stability conditions, and
transition rules. Without a comparative membrane, these systems cannot be safely
analyzed together.

\subsection{Related Work}

Several bodies of work inform the development of Shared-Horizon:

\begin{itemize}
    \item \textbf{Geometric AI and Operator Theory}  
    Research on holonomy, reversible operators, and geometric stability provides
    the foundation for NDH and Structural-AI systems.

    \item \textbf{Expressive Phenomenology}  
    ANIMA Classic (Zenodo DOI: 10.5281/zenodo.22812349) and the Anima GitHub
    repository offer expressive curvature models and phenomenological metrics
    that inform the expressive manifold definitions used in this preprint.

    \item \textbf{Narrative Dynamics}  
    Work on reversible narrative transitions, narrative coherence metrics, and
    expressive-narrative compatibility informs the RNE lineage.

    \item \textbf{Meta-System Analysis}  
    NDH Meta-Systems analysis
    (GitHub: Borealiscodes/NDH-Meta-Systems) provides mathematical baselines for
    operator geometry, expressive curvature, and lineage-safe comparison.

    \item \textbf{Comparative Frameworks}  
    Prior attempts at lineage comparison often rely on unifying abstractions,
    which risk collapsing lineage-specific invariants. Shared-Horizon instead
    preserves lineage sovereignty.
\end{itemize}

\subsection{Motivation for Shared-Horizon}

Existing comparative frameworks fail to account for:

\begin{enumerate}
    \item differences in curvature behavior,
    \item reversibility constraints,
    \item spectral drift boundaries,
    \item expressive geodesic structure,
    \item collapse-risk geometry,
    \item lineage-specific stability envelopes.
\end{enumerate}

Shared-Horizon addresses these gaps by introducing a neutral membrane that
preserves lineage invariants while enabling safe comparative analysis. The
framework is not a unification or synthesis; it is a structured comparative
surface that prevents conceptual collapse.

\subsection{Academic Use of External References}

All external references—including ANIMA Classic, the Anima GitHub repository,
and NDH Meta-Systems analysis—are used strictly for academic comparison. No code
or assets are incorporated, and no derivative work is produced. These references
provide mathematical baselines that ensure the correctness of lineage
comparisons within the Shared-Horizon framework.

\subsection{Background Summary}

The background establishes the fragmentation of AI lineages, reviews related
work, and motivates the need for a neutral comparative membrane. Shared-Horizon
builds on these foundations to provide a mathematically rigorous, lineage-safe
framework for analyzing diverse AI ecosystems.
